\section{Threats to Validity}
\label{sec:threats}

\paragraph{Construct validity.}
The taxonomy is an operational diagnostic for the prompt-and-validate
paradigm, not a complete account of structured-output failure. Primary
categories are single-valued, labels follow how a constraint is encoded
rather than what it intends (\cref{sec:setup:taxonomy}), and we quantify
co-occurrence only beneath fences. Classifier agreement is
development-set agreement from a single labeller who also wrote the
classifier. We evaluate conformance only, not whether extracted content
is correct. Among attempts whose model-level validators ran, no
cross-field violation occurred (\cref{app:crossfield}), but attempts
rejected earlier, including during JSON parsing, never reached those
checks, so the zero describes a filtered subset of the node-level
validators we implemented. Cross-node referential integrity, whether a
later node's identifiers appear in an earlier node's entities, is not
implemented in our schemas; it is expressible in Pydantic by passing the
earlier output as validation context
(\code{model\_validate\_json(data, context=\ldots)}) and is the natural
next measurement.

\paragraph{Internal validity.}
Repair attempts, which are conditioned on a prior failure, make up a
model-dependent share of each attempt rate; first-attempt rates and run
completion are reported separately for that reason
(\cref{app:breakdowns}). Every retry carries the validation error and no
condition retries without it, so the design does not separate the effect
of the feedback from that of a further attempt. Five models reject an
explicit temperature, and the Anthropic and OpenAI branches use
different token ceilings. Some questions were run more than once; keeping
only the earliest run per question changes any cell's failure rate by at
most \pn{robust.within_cell.max_abs_delta_pp} percentage points
(\cref{app:manifest}). The prompt-ablation variants were run after the
final baseline batch rather than interleaved with it and are not paired
by question.

\paragraph{External validity.}
The schemas are ours, and one ablation showed that a cap we set was
binding below the output lengths of one model; a different author would
produce a different distribution over categories, and the domain
comparison shows that a different task would too. We cover two vendors by
design, trading breadth for several generations and capability tiers per
vendor, which the within-family comparisons require; a third vendor or an
open-weights model may exhibit modes our taxonomy lacks. We do not
compare against constrained decoding or vendor structured-output modes
(\cref{sec:related}), so we make no claim about how often those
configurations would prevent the failures we report. Questions are fixed
dataset prefixes, the design is not fully crossed, and most failures in
this corpus come from a single model. Most endpoints returned an alias
rather than a pinned version, so later calls to the same identifier may
not reproduce these distributions.

\paragraph{Conclusion validity.}
Attempts are not independent, so all intervals are question-clustered
bootstrap intervals (\cref{sec:setup:basis}) that describe variability
over the observed question set, not over the benchmarks. No conclusion
depends on ordering models whose intervals overlap. We do not interpret
cells with a handful of failures, including the identifier-perturbation
cells and the combined length rates across prompt variants, and we treat
maximum-gap divergence statistics as secondary because they are selected
to be extreme (\cref{app:divergence}). The Haiku~4.5 FinanceBench cell
rests on \pn{c.claude-haiku-4-5@medium@financebench.questions} distinct
questions.
